## NHS2_itemsOfInterest_questions

## Section 1 of 6
On this form you'll help us to expand the Natural History of Song Discography to cover more cultures, languages, regions, and song types. You will fill out this form ONCE for each song you find.

You will decide to include recordings on the basis of the MATERIALS that are associated with that recording, NOT on the basis of how the recording sounds. In fact, you won't listen to the song until the very end of this process.

Stephanie, Connie, or Sam will assign you a region to focus on. Your job is to search a bunch of databases and libraries for recordings of several different behavioral functions (e.g., "for dancing"), that originate from that region. This form will help you to catalog each song along with some important information for each one.

Here are the basics

A.  What are your initials?

1. ABK
2. AGP
3. AHS
4. AIM
5. AMB
6. AMM
7. AVM
8. ANP
9. AZX
10. BSG
11. CMB
12. DIS
13. DMK
14. DPJ
15. DWL
16. ERA
17. EFK
18. HLR
19. HTD
20. JAY
21. JBA
22. JCS
23. JCVL
24. JDL
25. JMI
26. JNY
27. JSWB
28. KAH
29. KJH
30. KLL
31. KVL
32. LAV
33. LDR
34. LJD
35. LOT
36. LXN
37. MAK
38. MCB
39. MDS
40. MJW
41. MLM
42. MMX
43. NCM
44. NLC
45. NLR
46. NSB
47. PKH
48. PWC
49. RAH
50. REH
51. RJW
52. SAM
53. SAS
54. SAT
55. SDA
56. SFS
57. SHJ
58. SJG
59. SSG
60. TJS
61. TTI
62. WAU
63. WCS
64. XDF
65. YJS
66. YJW
67. YTD


B. What region are you searching in?
1. Amazon and Orinoco
2. Arctic and Subarctic
3. Australia
4. British Isles
5. Caribbean
6. Caucasus
7. Central Africa
8. Central America
9. Central Andes
10. Central Asia
11. Central Europe
12. Central Mexico
13. East Asia
14. Eastern Africa
15. Eastern Europe
16. Eastern South America
17. Eastern Woodlands
18. Maya Area
19. Melanesia
20. Micronesia
21. Middle East
22. North Asia
23. Northern Africa
24. Northern Mexico
25. Northwest Coast and California
26. Northwestern South America
27. Plains and Plateau
28. Polynesia
29. Scandinavia
30. South Asia
31. Southeast Asia
32. Southeastern Europe
33. Southern Africa
34. Southern Europe
35. Southern South America
36. Southwest and Basin
37. Western Africa
38. Western Europe

**How many songs should I catalog?**
You need to catalog at least THREE examples of each song type (e.g., "dance") from each region you are assigned and at most TEN examples of each song type. (Try your darndest to get up to 10 and don't just stop when you hit 3!)

If you can't find 3 examples in the published literature, then you'll need to work with Connie, Stephanie, Sam, Luke, and Manvir to contact anthropologists and ethnomusicologists who work in that region to see if any unpublished recordings are available. **Please do NOT contact anyone on your own! If you need to contact someone, STOP and ask for help.**

Sometimes you might find a lot of examples of a particular song type. If you do, to the extent possible, spread these examples across as many cultures as are available within your region. You should never catalog more than 10 examples of a given song type/region pairing.

## Section 2 of 6
Find a recording and tell us about it!
This page gives instructions for how to find a song to catalog, and then helps you to enter basic information about it.
Do NOT listen to the song yet!!
In NHS, we catalog songs based on the documentary information we have about them, NOT based on what you think of the song. You'll get a chance to listen to the song at the end of this form, but for now, do NOT listen to it!!! This is really important.

###Where to look for recordings
There are a few different ways that you'll be looking for recordings. Do them in the following order.
1. Search each of the GLOBAL music databases listed on Trello: https://trello.com/b/TivgB4cp
2. Search any relevant region- or culture-specific databases, also listed on Trello. (You should also see if you can find any other relevant databases (e.g., via Google searches) that might be useful).
3. Search libraries, starting with the WorldCat database (http://worldcat.org) and narrowing down to specific libraries to find items (including contacting librarians directly if you need additional information).
Once all of these are exhausted, we then start contacting anthropologists & ethnomusicologists directly to ask for recordings. If you get to that point, post to Slack to ask what to do.

###How to look for recordings
The region you're searching in is outlined visually in the map at https://goo.gl/xediiC.
a. If you are searching in a region in Central America, please see here for further clarification: https://drive.google.com/file/d/19GNCuNr8Jb6gQ8NGqghxrlp_ftZeMtrr/view?usp=sharing
b. If you are searching in a region in North America, please see here for further clarification: https://drive.google.com/file/d/17iFXOFEpb01RrgRx70Fjljxv3VgUvSq6/view?usp=sharing

The first way to search is simply to use a map from each music database where maps are available (e.g. Global Jukebox) and to look for examples from the region you're searching for. Easy!

**Note. When searching Global Jukebox, we now have the full metadata for each song available in a spreadsheet (see Trello). You may find it easier to search GJ by genre rather than relying solely on its map.**

If you can't find a map with songs already marked on it, however, you'll need to run text searches. Usually the easiest thing to do is search for a particular culture's name or a particular language group. There are a few ways to figure out what to search for.

1. You can find candidate cultures using the maps of D-PLACE (https://d-place.org/ → choose “PLACE” → select region(s) of interest and click SEARCH). Find cultures in that region and then search the resources listed above for tracks.

### **important note!**
Many cultures have a large variety of spellings of their names. This makes it hard to find stuff. You should also search for "ethonyms" — alternative spellings of culture names — by using Wikipedia as well as by browsing the HRAF culture listing, available at http://ehrafworldcultures.yale.edu/ehrafe/browseCultures.do?context=main

2. You can do the same thing with languages: look for a particular language by searching D-PLACE for different languages that tend to fall in a particular region (which you can check out using D-PLACE's map feature). Then, search the resources above for recordings that have that language as a keyword.

3. If you find a track associated with a culture but are unsure whether it falls into the corresponding region, you can search that culture in D-PLACE. If the culture is not in D-PLACE, search the culture in Ethnologue to find an ISO code for its language, search Glottolog for that ISO code to get rough latitude & longitude, plug that into Google Maps and compare to the HRAF map to get its region. (If you can't find information at all about a given culture, ask Luke & Manvir for help).

### Important things
No tracks should be catalogued more than once. Before you catalog a song, make sure it isn't already in NHS! You can do that by checking the full listing of NHS tracks at https://goo.gl/TFm79g.

If you have exhausted all available recordings for a particular region, and have not found 3 examples of each song type, then we'll need to contact anthropologists and ethnomusicologists to find unpublished recordings. Don't do this yourself. Post to Slack to find out what to do.

Here are the songs you're looking for:
There are now TEN categories! Sometimes songs cover more than one category: in those cases, your job will be to decide which category fits the song *best*. Put that category as the primary function, and any other category as a secondary function.

If you really can't be sure whether a song qualifies for a particular primary category, you can still catalog it, but specify your concern in the notes field at the end of the form.

DANCE
Inclusion criteria: Sung with goal of a person or persons dancing along to it.
Exclusion criteria: Songs that happen to be accompanied by dancing but are used for other goals.

HEALING
Inclusion criteria: Sung in a healing ceremony with the goal of curing sickness.
Exclusion criteria: Songs describing sick people or a past epidemic.

LOVE
Inclusion criteria: Sung to express love directly to another person or to describe currently felt love.
Exclusion criteria: Songs about unrequited love, deceased loved ones, or love for animals or property

LULLABY
Inclusion criteria: Sung to an infant or child with the goal of soothing, calming, or putting to sleep.
Exclusion criteria: Songs designed to excite the listener (e.g., “play songs”); singing games

PLAY
Inclusion criteria: A song that excites a child or infant and and engages them in play. This can include singing games.
Exclusion criteria: Children's songs for soothing, calming, or putting to sleep.

PROCESSION
Inclusion criteria: sung to accompany a formalized march, entrance, or parade, such as during a wedding, funeral or the introduction of a leader [*note*: please specify in 'secondary function' the type of event that the procession is associated with, e.g., "wedding"]
Exclusion criteria: Processions of dancing

MOURNING
Inclusion criteria: sung to express grief or sadness about the death of a person, in the present or past.
Exclusion criteria: Songs for sick or dying people (these are more likely to be healing songs), or laments about events other than the death of a person.

WORK
Inclusion criteria: sung to accompany work activities, including planting, grinding, harvesting, processing, tool-making
Exclusion criteria: Hunting songs (esp. songs sung to celebrate successful hunts, songs sung to prepare for hunts)

STORY
Inclusion criteria: sung to recount historical or mythological events, narrate a sequence of activities by one or more persons, etc. (e.g., a ballad)
Exclusion criteria: Lullabies that include stories

PRAISE
Inclusion criteria: sung to express admiration for the traits or accomplishments of a person, animal, location, or item of property
Exclusion criteria: A song expressing love for another person or explicitly religious songs (like devotionals).

When you are searching for songs, make your best guess as to whether or not a song fits one of these categories. In general, it's better to include a song that you're unsure of rather than not include it: other people will be reviewing your decisions later and trimming down the dataset to only songs that we're clear on, so it's ok to include songs that you aren't 100\% sure on.

OK, now you're ready to catalog!
If you have any questions about the above information, STOP and ask before cataloguing.

If you're ready to roll, enter the info for your song below:

C. What type of song is this?
1. Dance
2. Healing
3. Love
4. Lullaby
5. Play
6. Procession
7. Mourning
8. Work
9. Story
10. Praise

D. Is there any secondary function for this song?
1. No secondary function
2. Dance
3. Healing
4. Love
5. Lullaby
6. Play
7. Procession
8. Mourning
9. Work
10. Story
11. Praise

E. What source(s) is this song listed in ? (If it's in multiple places, you can check more than one box)
1. None of these
2. Global Jukebox
3. Garland Encyclopedia of World Music
4. Smithsonian Folkways
5. CREM
6. British Library
7. WorldCat

F. Is this an unpublished track that came directly from an anthropologist?
1. Yes
2. No

## Section 3 of 6
Who is the anthropologist/ethnomusicologist?

G. Name
[free field]

H. Affiliation
[free field]

I. Email
[free field]

## Section 4 of 6
How do you know that this song is what you think it is?
Some recordings have an informative label: someone else has categorized the song and has stated this category, e.g., "love song".
Some recordings have supporting text (ethnography).
Some recordings have lyrics, which you can interpret (e.g., "I love you! I love you! I love you!" is probably a love song). Some songs have a combination of these things.

This page helps us to understand why you categorized the song in the way that you did.

J. Tell us about how you categorized this song. (check all that apply)
1. The song is identified as such in its title or label
2. There is text in associated ethnography that states what type of song this is (this could be in liner notes, it could be in an email from an anthropologist, etc)
3. The lyrics made it clear to me what type of song this is

Supporting info
At the end of your entry, specify the DETAILED SOURCE of this label and enclose that information in //CAPITAL LETTERS WITHIN DOUBLE SLASHES//.

For example, if the song included supporting ethnography, you might type:
The singer of this song performs it when a child is sick so that the child will feel better. //CD HAS ADDITIONAL LINER NOTES BOOKLET, THIS IS ON PAGE 4//

When specifying the DETAILED SOURCE, you should NOT simply repeat the name of the database you got this from (e.g., //GLOBAL JUKEBOX//). Rather, give something that will make it very easy for us to find the information you coded (e.g., //CODING ID 3412 IN GLOBAL JUKEBOX, RELEVANT TEXT IS IN THE DESCRIPTION FIELD//)

K. Enter the label of the song here.
If the song was labeled in the source (most are), enter the exact label of the song here (e.g., "Lullaby"). Sometimes this is just the name of the track, but other times, a database might have a field for this. If a label is not available, enter "NA". ****Don't forget to include the //SOURCE// for this information.**** [Note. The song title is not necessarily a label. A label is informative: "Golden Slumbers" is not a label. "Lullaby" is a label. Please do not enter a label if it is not informative.]
[free field]

L. Information from ethnography
If there was associated ethnographic text informing your decision about the song type, copy and paste that text here. If supporting ethnographic text is not available, enter "NA".  ****Don't forget to include the //SOURCE// for this information.****
[free field]

M. Information from lyrics
If you used the lyrics to inform your decision about the song type, copy and paste the lyrics here. If lyrics are not available, enter "NA". ****Don't forget to include the //SOURCE// for this information.****
[free field]

## Section 5 of 6
Metadata
This section gives general information about this song so that we can easily find it later.

N. Paste a hyperlink here that gets to the source material (e.g., the CD containing the track).
Make *sure* that this hyperlink works: it should lead directly to information about the source material (e.g., a WorldCat permalink to the album from which this song was drawn) and NOT just to the database that the material was collected from (e.g., a general link to Global Jukebox). For database searches, do your best to find a direct link to the song. For items sent directly to you by anthropologists/ethnomusicologists, upload it to the Music Lab NAS (/musiclabnas/NHS2/submissions) and generate a link for it.
[free field]

O. If necessary, provide additional information here about how to find this track.
Only use this field if it's useful/important, e.g., "The WorldCat link refers to a book; inside that book there is an attached CD that contains the track." Do not put anthropologists' contact info here.
[free field]

P. List the track number of the song.
Most songs that come from a CD or album will require this, e.g., "Track 4". If you need to specify additional clarifying info, please do ("Disc 2, Track 4"). If there isn't a track number (i.e., the permalink resolves to this song and nothing else) then enter "NA"
[free field]

Q. Enter the track name, as it appears in the source.
This is usually the actual song title (e.g., "Mana gallaka noukat") but sometimes is the same as the label (e.g., "Lullaby")
[free field]

R. Enter the name of the culture this song was collected from.
Use the information provided by the author of the source material. If there are multiple culture names, or if you're not sure, give as much information here as you can.
[free field]

S. General notes
Please use this field sparingly. You can write anything that is important to know about this track here. For instance: "Global Jukebox says the audio is available, but it doesn't play properly online. We might need to find the original recording of this one."
[free field]

T. Do you have access to the audio of this song?
If you can listen to the song (whether by listening to an online recording, or waltzing over to a library to check it out) then click "Yes" to answer a few questions about it.
1. Yes (--> Go to section 6 (Audio info))
2. No (--> Submit form)

### Section 6 of 6
Audio info
Nearly there — this is the last page.

U. Confirm here that you did not listen to the song until this point!
If you listened to the song before this point, even though you were not supposed to, explain why in the Other field.
1. I followed the instruction and I have not yet listened to it!

Yay, you can finally listen to the song!
You don't need to listen to the whole thing if you don't want to; just listen enough to answer the two questions below.

V. Is there at least one audible singer on the recording?
You can interpret "singer" fairly liberally; e.g., if there's an audible person but it sounds more like "chanting" than "singing", you can still choose "Yes".
1. Yes
2. No

W. How easy is it to hear the singer(s)?  
1 (Very difficult); 2; 3; 4; 5; 6 (Very easy)
